\documentclass[../CLT_SYK.tex]{subfiles}
\setcounter{section}{2}

\begin{document}

\section{Direct Rademacher control of all higher chaoses} \label{sec:Rademacher-remainder}

Following the expressions of the first and second terms in in the polynomial chaos expansion \eqref{eq: chaos} of $Z_N(\beta)$,
we estimate the higher-order terms in this section.
Recall definition of $Y_N(\beta)$ in \eqref{eq:Y_N} and the orthogonal projection $P_k$ defined in Section~\ref{sec:chaos-decomposition}.

\begin{proposition} \label{prop:higher-chaos-closure}
  There exists $\beta_{\mathrm{r}}=\beta_{\mathrm{r}}(q)>0$ such that, for every fixed real $\beta$ with $|\beta|<\beta_{\mathrm{r}}$,
  \begin{equation*}
    \sum_{k \geq 3}
    \E\mleft[\left|P_k Y_N(\beta)\right|^2\mright]
    =
      O\mleft(N^{-3(q/2-1)}\mright)
    =
      o\mleft(N^{2-q}\mright).
  \end{equation*}
\end{proposition}

To prove this proposition, we use the following two lemmas.

\begin{lemma} \label{lem:Z-derivative-bound}
  For any $\beta\in \mathbb R$ and for any multi-index $\alpha$, it holds 
  \begin{equation}\label{eq:Z-derivative-bound}
    \left|\partial_G^\alpha Z_N(\beta;G)\right|
    \leq (|\beta|\sigma_N)^{|\alpha|} Z_N(\beta;G).
  \end{equation}
\end{lemma}

\begin{proof}
  Let $k = |\alpha|$.
  The case that $k=0$ is trivial, hence assume that $k \geq 1$ in the following.
  We claim that, for any $(I_i)_{i=1}^{k} \subset \mathcal I_N$,
  \begin{multline}
    \partial_{g_{I_k}}\cdots\partial_{g_{I_1}}Z_N(\beta;G) \\
    =
      (-\beta\sigma_N)^k
      \sum_{\pi\in \mathfrak{S}_k}\int_{\Delta_k}
      \tr\mleft(
        e^{-t_0\beta H_N}
        \psi_{I_{\pi(1)}}
        e^{-t_1\beta H_N}
        \cdots
        \psi_{I_{\pi(k)}}
        e^{-t_k\beta H_N}
      \mright)
      \odif{t_1}\cdots\odif{t_k},
    \label{eq:iterated-Duhamel-formula}
  \end{multline}
  where
  \[
    \Delta_k =
    \left\{
      (t_j)_{j=1}^{k} \in \interval{0}{1}^k
      :
      \sum_{j=1}^{k} t_j\le 1
    \right\},
    \qquad
    t_0 = 1 - \sum_{j=1}^{k} t_j.
  \]
  By the equality that $\partial_{g_I}H_N = \sigma_N\psi_I$ and the Duhamel formula~\cite[Corollary~3.11]{HiaiPetz2014}, for any $I \in \mathcal I_N$ and any $t \geq 0$,
  \begin{align}
    \partial_{g_I}e^{-t\beta H_N}
    &=
      -t\beta\sigma_N\int_0^1
      e^{-st\beta H_N}\psi_I
      e^{-(1-s)t\beta H_N}\,\odif{s} \notag\\
    &=
      -\beta\sigma_N\int_0^t
      e^{-u\beta H_N}\psi_I
      e^{-(t-u)\beta H_N}\,\odif{u}.
    \label{eq:Duhamel-interval-splitting}
  \end{align}
  Note that  \eqref{eq:Duhamel-interval-splitting} with $t=1$ yields
  \eqref{eq:iterated-Duhamel-formula} for $k=1$.
  Suppose inductively that the formula holds at order $k - 1$.
  Then
  \begin{multline}
    \partial_{g_{I_k}}\cdots\partial_{g_{I_1}}
    Z_N(\beta;G)
    =
      (-\beta\sigma_N)^{k-1}
      \sum_{\pi\in \mathfrak{S}_{k-1}}
      \int_{\Delta_{k-1}} \odif{t_1}\cdots\odif{t_{k-1}} \\
      \times
      \partial_{g_{I_k}}
      \tr\Bigl(
        e^{-t_0\beta H_N}
        \psi_{I_{\pi(1)}}
        e^{-t_1\beta H_N}
        \cdots
        \psi_{I_{\pi(k-1)}}e^{-t_{k-1}\beta H_N}
      \Bigr)
    \label{eq:iterated-Duhamel-formula-induction}
  \end{multline}
  Fix $\pi \in \mathfrak{S}_{k-1}$ and $j\in\{1, \ldots, k\}$.
  Differentiating the factor $e^{-t_{j-1}\beta H_N}$,
  and relabeling the running variables,
  the Leibniz rule leads to
  \begin{align*}
    -\beta\sigma_N
    \int_{\Delta_k}
    \tr\mleft(
      e^{-t_0\beta H_N}
      \psi_{I_{\tilde{\pi}_{j}(1)}}
      e^{-t_1\beta H_N}
      \cdots
      \psi_{I_{\tilde{\pi}_{j}(k)}}
      e^{-t_k\beta H_N}
    \mright)
    \odif{t_1} \cdots \odif{t_k},
  \end{align*}
  where
  \[
    \tilde{\pi}_{j}(k)
    =
    \begin{cases}
      \pi(k),   & k<j, \\
      k,        & k=j, \\
      \pi(k-1), & k>j.
    \end{cases}
  \]
  Since $(\pi, j) \mapsto \tilde{\pi}_{j}$ is a bijection from
  $\mathfrak{S}_{k-1}\times\{1,\ldots,k\}$ onto $\mathfrak{S}_k$, summing the terms produced by the Leibniz rule in
  \eqref{eq:iterated-Duhamel-formula-induction} proves
  \eqref{eq:iterated-Duhamel-formula}.

  To bound the integrand in \eqref{eq:iterated-Duhamel-formula}, use the normalized Schatten norms for a matrix $A$ and $p \in \interval[open right]{1}{\infty}$:
  \[
    \norm{A}_p \coloneqq \left(\tr\mleft(\sqrt{A^* A}^p\mright)\right)^{1/p}.
  \]
  Note that the case that $p=\infty$ corresponds to the operator norm $\norm{A}_\infty$.
  Let $t_0, \ldots, t_k \in \interval{0}{1}$ satisfying $\sum_{j=0}^k t_j = 1$.
  Iterating H\"older's inequality~\cite[Theorem~6.20]{HiaiPetz2014} for any matrices $A_0, \ldots, A_k$ and $B_1, \ldots, B_k$ with the finiteness of the Schatten norm provides
  \begin{align*}
    \left|\tr(A_0 B_1 A_1 \cdots B_k A_k)\right|
    & \leq
      \norm{A_0 B_1 A_1 \cdots B_k A_k}_1\\
    & \leq
      \norm{A_0 B_1 A_1 \cdots A_{k-1} B_k}_{1/(t_0+\cdots+t_{k-1})}
      \norm{A_k}_{1/t_k} \\
    & \leq
      \norm{A_0 B_1 A_1 \cdots B_{k-1} A_{k-1}}_{1/(t_0+\cdots+t_{k-1})}
      \norm{B_k}_\infty
      \norm{A_k}_{1/t_k} \\
    & \leq
      \cdots \\
    & \leq
      \prod_{j=0}^{k} \norm{A_j}_{1/t_j}
      \prod_{i=1}^{k} \norm{B_i}_\infty.
  \end{align*}

  We apply this inequality with $A_j = e^{-t_j \beta H_N}$ and $B_i = \psi_{I_{\pi(i)}}$.
  Note that $H_N$ is Hermitian and $\norm{\psi_I}_\infty = 1$ due to $\psi_I^*\psi_I =\mathbb I$. 
  Thus, for every $\pi \in \mathfrak{S}_k$ and $(t_j)_{j=1}^{k} \in \Delta_k$,
  \begin{align}
    \MoveEqLeft
    \left|
      \tr\mleft(
        e^{-t_0\beta H_N}
        \psi_{I_{\pi(1)}}
        e^{-t_1\beta H_N}
        \cdots
        \psi_{I_{\pi(k)}}
        e^{-t_k\beta H_N}
      \mright)
    \right| \notag \\
    \leq {} &
      \prod_{j=0}^{k}
      \norm{e^{-t_j\beta H_N}}_{1/t_j}
      \prod_{i=1}^{k}
      \norm{\psi_{I_{\pi(i)}}}_\infty \notag\\
    \leq {} &
      \prod_{j=0}^{k}
      \tr\mleft(e^{-\beta H_N}\mright)^{t_j}
    =
      \tr\mleft(e^{-\beta H_N}\mright)
    =
      Z_N(\beta;G),
    \label{eq:Duhamel-integrand-Holder-bound}
  \end{align}
  where we used $\sum_{j=0}^{k} t_j=1$.
  Combining \eqref{eq:iterated-Duhamel-formula} and \eqref{eq:Duhamel-integrand-Holder-bound}, we obtain
  \begin{align*}
  \left|
    \partial_{g_{I_k}}
    \cdots
    \partial_{g_{I_1}}
    Z_N(\beta;G)
  \right|
  \leq {} &
    \left(|\beta| \sigma_N\right)^k
    \sum_{\pi\in \mathfrak{S}_k}
    \int_{\Delta_k}
    Z_N(\beta;G)
    \odif{t_1} \cdots \odif{t_k} \\
  = {} &
    \left(|\beta| \sigma_N\right)^k
    Z_N(\beta;G)
    |\mathfrak{S}_k|
    \int_{\Delta_k} 1 \odif{t_1} \cdots \odif{t_k} \\
  = {} &
    \left(|\beta| \sigma_N\right)^k
    Z_N(\beta;G),
  \end{align*}
  which proves \eqref{eq:Z-derivative-bound}.
\end{proof}

Recall that Lemma~\ref{lem: parity rule lemma} shows that only even multi-indices have nonzero Hermite coefficients.
Let $(\varepsilon_x)_{x=1}^N$ be independent Rademacher variables on an auxiliary probability space.
Write $\mathbb{P}_\varepsilon$ and $\E_\varepsilon$ for the probability measure on this space and the associated expectation, respectively.
Set
\begin{equation}\label{eq:Rademacher-chaos-definition}
  S_{N,q}
  \coloneqq
    \sum_{I\in\calI_N}\varepsilon_I,
  \qquad
  \varepsilon_I
  \coloneqq
    \prod_{x\in I}\varepsilon_x.
\end{equation}

\begin{lemma} \label{lem:Rademacher-encoding}
  For every integer $k \geq 0$,
  \begin{equation}
    \label{eq:even-sum-Rademacher-moment}
    \sum_{
      \substack{
        |\alpha|=k \\
        \alpha\ \mathrm{even}
      }
    }
    \frac{1}{\alpha!}
    =
      \frac{1}{k!}
      \E_\varepsilon\mleft[S_{N,q}^k\mright].
  \end{equation}
  Moreover, there are constants $c_q > 0$ and $C_q < \infty$ depending only on $q$ such that, for every $N$ and every $0 \leq t \leq c_q N^{-(q-1)}$,
  \begin{equation}
    \label{eq:Rademacher-exponential-moment}
    \E_\varepsilon\mleft[
      |S_{N,q}|^3
      e^{t |S_{N,q}|}
    \mright]
    \leq
      C_q N^{3q/2}.
  \end{equation}
\end{lemma}

\begin{proof}
  We first show \eqref{eq:even-sum-Rademacher-moment}.
  Starting from the definition \eqref{eq:Rademacher-chaos-definition}
  and applying the multinomial formula provides
  \begin{equation}\label{eq:Rademacher-multinomial-expansion}
    \frac{1}{k!}\E_\varepsilon\mleft[S_{N,q}^k\mright]
    =\frac{1}{k!}\E_\varepsilon\mleft[
      \sum_{|\alpha|=k}k!
      \prod_{I\in\calI_N}\frac{\varepsilon_I^{\alpha_I}}{\alpha_I!}
    \mright]
    =
    \sum_{|\alpha|=k}\frac{1}{\alpha!}
    \E_\varepsilon\mleft[
      \prod_{I\in\calI_N}\varepsilon_I^{\alpha_I}
    \mright].
  \end{equation}
  For a fixed $\alpha$, the product inside the expectation is
  \begin{equation*}
    \prod_{x=1}^N
    \varepsilon_x^{\sum_{I: x\in I}\alpha_I}.
  \end{equation*}
  By independence of the Rademacher variables and the definition of an even
  multi-index in \eqref{eq:even-multi-index}, its expectation is one if
  $\alpha$ is even and zero otherwise.  Substituting this into
  \eqref{eq:Rademacher-multinomial-expansion} proves
  \eqref{eq:even-sum-Rademacher-moment}.

  We next bound $S_{N,q}$.  For every integer $N\geq1$ and $m\geq0$, let $e_m^{(N)}$ denote the $m$th elementary symmetric polynomial in $N$ variables:
  \begin{equation*}
    e_m^{(N)}(z_1,\ldots,z_N)
    \coloneqq
      \sum_{
        \substack{
          I\subset[N]\\
          |I|=m
        }
      }
      \prod_{x\in I}
      z_x,
    \qquad
    e_0^{(N)}(z_1,\ldots,z_N)
    \coloneqq
      1.
  \end{equation*}
  For the rest of the proof, we abbreviate $e_m \coloneqq e_m^{(N)}(\varepsilon_1,\ldots,\varepsilon_N)$.
  Then, $e_1=\sum_{x=1}^N\varepsilon_x$ and \eqref{eq:Rademacher-chaos-definition} imply that $S_{N,q}=e_q$.
  Newton's identities (see, e.g., \cite[Chapter I, \S2]{Macdonald1995}), state that
  \begin{equation}\label{eq:Newton-identities}
    m e_m
    =\sum_{j=1}^m(-1)^{j-1}e_{m-j}p_j,
  \end{equation}
  where $p_j \coloneqq \sum_{x=1}^N\varepsilon_x^j$.  Since
  $\varepsilon_x\in\{\pm1\}$, we have
  \begin{equation*}
    p_j
    =
      \begin{cases}
        e_1, & j\text{ odd},  \\
        N,   & j\text{ even}.
      \end{cases}
  \end{equation*}
  We now prove by induction on $m$ that there is a constant $C_m < \infty$ such that
  \begin{equation*}
    |e_m|
    \le C_m\left(N^{m/2}+|e_1|^m\right).
  \end{equation*}
  Suppose that the estimate holds for  orders less than $m$.
  Set $X_N \coloneqq e_1 / \sqrt N$.
  For every integer $j\ge1$,
  \[
    |p_j| \le N^{j/2} \left(1 + |X_N|^j\right),
  \]
  hence, by the induction hypothesis and putting $C_0 = 1$,
  \begin{align*}
    |e_{m-j}p_j|
    &\le
      C_{m-j}N^{m/2}
      (1+|X_N|^{m-j})(1+|X_N|^j)\\
    &\le
      2C_{m-j}N^{m/2}(1+|X_N|^m),
  \end{align*}
  where we used $1+|X|^m-|X|^{m-j}-|X|^j = (1-|X|^{m-j})(1-|X|^j) \ge 0$.
  Substituting this into \eqref{eq:Newton-identities} proves the induction step.
  Thus, taking $m=q$ and using \eqref{eq:Rademacher-chaos-definition}, we obtain
  \begin{equation}\label{eq:Rademacher-chaos-deterministic-bound}
    |S_{N,q}|
    \leq C_q\left(N^{q/2}+|e_1|^q\right).
  \end{equation}

  Since $q\geq4$ and $|X_N|\leq\sqrt N$,
  \eqref{eq:Rademacher-chaos-deterministic-bound} implies that, for $0\leq t\leq c_qN^{-(q-1)}$,
  \begin{equation}
    t|S_{N,q}|
    \leq
      C_q c_q N^{-(q/2-1)} \left(1 + |X_N|^q\right)
    \leq
      C_q c_q \left(1 + |X_N|^2\right).
    \label{eq:Rademacher-exponent-reduction}
  \end{equation}
  Indeed,
  \begin{equation*}
    N^{-(q/2-1)}|X_N|^q
    =|X_N|^2
    \left(\frac{|X_N|}{\sqrt N}\right)^{q-2}
    \leq |X_N|^2.
  \end{equation*}
  Choose $c_q$ small enough so that $C_q c_q\leq 1/4$.  Hoeffding's inequality~\cite[Theorem~2]{Hoeffding1963} provides
  \begin{equation}\label{eq:Hoeffding-Rademacher-sum}
    \mathbb P_\varepsilon(|X_N|\geq u)
    \leq
      2e^{-u^2/2}
    \qquad
    (u\geq0).
  \end{equation}
  Consequently, integrating the tail bound in \eqref{eq:Hoeffding-Rademacher-sum} yields
  \begin{equation}\label{eq:uniform-Rademacher-integrability}
    \sup_{N\geq1}
    \E_\varepsilon\mleft[
      \left(1+|X_N|^{3q}\right)
      e^{|X_N|^2/4}
    \mright]
    <
    \infty.
  \end{equation}
  Finally, \eqref{eq:Rademacher-chaos-deterministic-bound}, \eqref{eq:Rademacher-exponent-reduction} and \eqref{eq:uniform-Rademacher-integrability} imply that
  \begin{align*}
    \E_\varepsilon\mleft[
      |S_{N,q}|^3e^{t|S_{N,q}|}
    \mright]
    &\leq
    C_qN^{3q/2}
    \E_\varepsilon\mleft[
      \left(1+|X_N|^{3q}\right)
      e^{(1+|X_N|^2)/4}
    \mright] \\
    &\leq C_qN^{3q/2},
  \end{align*}
  which proves \eqref{eq:Rademacher-exponential-moment}.
\end{proof}

\begin{proof}[Proof of Proposition~\ref{prop:higher-chaos-closure}]
  For any $k = |\alpha| \geq 3$, the polynomial chaos expansion \eqref{eq: chaos} and its orthogonality relation lead to
  \[
    \E\mleft[\left|P_k Y_N(\beta)\right|^2\mright]
    =
      \frac{1}{\E[Z_N(\beta)]^2}
      \sum_{|\alpha|=k}\frac{1}{\alpha!}
      \left|
        \E\mleft[Z_N(\beta) H_\alpha(G)\mright]
      \right|^2.
  \]
  By Lemmas~\ref{lem: parity rule lemma} and~\ref{lem:Gaussian-Hermite-IBP},
  \[
    \E\mleft[Z_N(\beta)H_\alpha(G)\mright]
    =
      \begin{cases}
        0,                                & \alpha \text{ is odd},  \\
        \E[\partial_G^\alpha Z_N(\beta)], & \alpha \text{ is even}.
      \end{cases}
  \]
  Moreover, Lemma~\ref{lem:Z-derivative-bound} provides,
  \[
    \bigl|\E[\partial_G^\alpha Z_N(\beta)]\bigr|
    \leq
      \E\bigl[\left|\partial_G^\alpha Z_N(\beta)\right|\bigr]
    \leq
      (|\beta|\sigma_N)^{k}
      \E\mleft[Z_N(\beta)\mright].
  \]
  Since $\E[Z_N(\beta)] > 0$, it follows that
  \begin{equation}
    \E\mleft[\left|P_k Y_N(\beta)\right|^2\mright]
    =
      \frac{1}{\E[Z_N(\beta)]^2}
      \sum_{
        \substack{
          |\alpha| = k \\
          \alpha\ \mathrm{even}
        }
      }
      \frac{1}{\alpha!}
      \bigl|\E\mleft[\partial_G^\alpha Z_N(\beta)\mright]\bigr|^2
    \leq (|\beta| \sigma_N)^{2 k}
      \sum_{
        \substack{
          |\alpha| = k \\
          \alpha\ \mathrm{even}
        }
      }
      \frac{1}{\alpha!}.
    \label{eq:Hermite-Rademacher-bound}
  \end{equation}
  By Lemma~\ref{lem:Rademacher-encoding}, summing \eqref{eq:Hermite-Rademacher-bound} over all $k \geq 3$ yields
  \begin{align}
    \sum_{k\geq3}\E\mleft[\left|P_k Y_N(\beta)\right|^2\mright]
    & \leq
      \sum_{k\geq3}
      \frac{(|\beta| \sigma_N)^{2 k}}{k!}
      \E_\varepsilon\mleft[S_{N,q}^k\mright] \notag\\
    & \leq
      \E_\varepsilon\mleft[
        \sum_{k\geq3}
        \frac{(\beta^2 \sigma_N^2 |S_{N,q}|)^k}{k!}
      \mright] \notag\\
    & \leq
      \frac{(|\beta| \sigma_N)^{6}}{6}
      \E_\varepsilon\mleft[
        |S_{N,q}|^3
        e^{\beta^2 \sigma_N^2 |S_{N,q}|}
      \mright],
    \label{eq:all-higher-chaoses-at-once}
  \end{align}
  where the second inequality uses $|\E_\varepsilon[S_{N,q}^k]| \leq \E_\varepsilon[|S_{N,q}|^k]$,
  and the last inequality uses $\sum_{k\geq3}u^k / k! \leq u^3e^u/6$ for $u\geq0$.
  Choose
  \begin{equation*}
    \beta_{\mathrm{r}}
    \coloneqq
      \sqrt{\frac{c_q}{(q-1)!}},
  \end{equation*}
  where $c_q$ is the constant in Lemma~\ref{lem:Rademacher-encoding}.
  If $|\beta| < \beta_{\mathrm{r}}$, then $\beta^2\sigma_N^2 = \beta^2(q-1)!N^{-(q-1)} \leq c_q N^{-(q-1)}$.
  Therefore the second claim of Lemma~\ref{lem:Rademacher-encoding}, together with \eqref{eq:all-higher-chaoses-at-once}, implies
  \begin{align*}
    \sum_{k\geq3}
    \E\mleft[
      \left|P_k Y_N(\beta)\right|^2
    \mright]
    &\leq
      C_{q, \beta}
      N^{-3(q-1)}
      N^{3q/2}
    =
      C_{q, \beta}
      N^{-3(q/2-1)}.
  \end{align*}
  Since $q \geq 4$, we have $3 (q / 2 - 1) > q - 2$, which completes the proof.
\end{proof}

\end{document}
